\section{はじめに}
\label{sec:intro}

物体を一度に一つの光源で撮影する one light at a time（OLAT）方式と、新たな視点からの再照明は、写実的なデジタルアセットを作成するための最も古くからある方法の一つである~\cite{debevec2000acquiring}。
近年では、各画像を一つの校正済み点光源で照らした数百から数千枚の画像から再照明可能な表現を再構成し~\cite{bi2020neural,zeng2023nrhints}、Gaussian Splatting~\cite{kerbl20233d} によってその結果を実時間でレンダリングできるようになった~\cite{bi2024gs3,kuang2024olat,fan2025rng,liu2025bigs,zheng2026ssdgs}。

点光源下では、画像が自然に見えるかどうかを左右する要因の多くは、材質を評価する前に決まる。すなわち、光がどこに当たり、どこに当たらないかである。
現在のガウシアンを用いた手法は、これをプリミティブごとに決定する。
\gsthree~\cite{bi2024gs3} と SSD-GS~\cite{zheng2026ssdgs} は、光源に向けてシーンをスプラットし、各ガウシアンに一つの可視性値を与える。RNG~\cite{fan2025rng} は深度シャドウマップを一つのスカラー手掛かりに集約し、半透明性は局所的な入力からガウシアンごとのネットワークで予測される~\cite{dihlmann2024sss,zheng2026ssdgs}。
そのため、床上の大きなガウシアンは一つの影の値しか持てず、移動する影の境界を表現するには、その都度高密度化によって形状を細分する必要がある。

本研究では、光源の画像平面上で考える。
点光源からガウシアンをラスタライズすると、アルファ合成重みは、各光源画素の光束をそのレイ上にあるプリミティブに分配する。スプラッティングモデルの範囲内では、これらの重みは各ガウシアンが遮り受け取る光束の割合である。
したがって、一回の追加ラスタライズにより、各プリミティブがどれだけの光をどの深度で吸収したかが記録される。本研究ではこれを光束の\emph{沈着}（deposit）と呼ぶ。
順方向レンダリングでは、既知の幾何形状上に手設計の量を記録するシャドウマップ~\cite{williams1978casting,donnelly2006variance,peters2015moment}、translucent shadow map~\cite{dachsbacher2003translucent}、reflective shadow map~\cite{dachsbacher2005reflective} において、数十年にわたり光源視点がこのように使われてきた。本研究ではこれを、幾何形状と記録する量をともに学習する逆レンダリングへ導入する。

提案手法 \ours（Light-Space Atlas）は、光源ごとに沈着の一次・二次深度モーメントと少数の学習した光束特徴をスプラットし、フィルタリングしてピラミッドを構築する（\cref{fig:teaser}）。
シェーディングを評価する任意の点、すなわち\emph{受光点}（receiver）は、光源視点に投影され、このピラミッドを参照する。幾何形状の最適化中はガウシアン中心を、その後は画素を受光点とする。
深度モーメントからは variance shadow map~\cite{donnelly2006variance} に倣ったフィルタ付きシャドウ判定を得て、残差ネットワークで補正する。光束特徴は沈着した特徴に対して線形な光輸送項に入力され、これは translucent shadow map の学習版に相当する。
アトラスには画像ごとに一回の $512^2$ ラスタライズが必要であり、入力は光源に対する相対量で記述され、勾配はアトラスを通じてガウシアンまで伝播する。

表示用に符号化した放射輝度上の損失でこのモデルを学習すると、Lego のブルドーザーのバケットのような細く入り組んだ構造が塊状に崩れた（\cref{fig:loss_domain}）。
そこで、放射量に基づいて最適化を段階化する。まず、ネットワークを用いない点光源モデルで幾何形状を当てはめ、その間に目標値を表示用放射輝度から線形放射輝度へ徐々に移行する。これは \gsthree と SSD-GS が HDR データに対して行う処理に相当する~\cite{bi2024gs3,zheng2026ssdgs}。次に、ガウシアン受光点を用いて完全モデルを線形放射輝度上で学習する。最後に、幾何形状を固定し、画素受光点で外観を微調整する。
二つのシーンの対照実験では、損失を計算する領域だけを変えることで、ホールドアウトしたビューの PSNR が Lego で 23.0 から 30.2\,dB、Drums で 24.1 から 30.1\,dB に向上した。

三つの公開 OLAT ベンチマーク~\cite{zeng2023nrhints,bi2024gs3,dihlmann2024sss} の全 18 シーンにおいて、\gsthree、SSD-GS、OLAT Gaussians~\cite{kuang2024olat}、RNG を共通の評価手順で再学習し、同一の目標画像で評価したところ、\ours は平均 PSNR、SSIM、LPIPS のすべてで最良の結果を達成した（PSNR は 33.47\,dB、次点は 31.59\,dB）。学習時間はシーンあたり約 20 分である。
ベンチマーク別では、最良のベースラインに対して実撮影データで 2.0\,dB（ビューごとの 2D 位置合わせ後は 0.7\,dB）、\gsthree シーンで 0.4\,dB 上回る一方、SSS-GS では OLAT Gaussians を 0.1\,dB 下回る。このため、すべてのベンチマークで上位二位以内に入る唯一の手法となる。
アブレーション実験は光輸送項の限界も示す。半透明材質では特に重要であるが、ほとんどの他の材質では同規模の局所ネットワークが同等の結果を得る。
本研究の貢献は以下のとおりである。
\begin{itemize}
\item ガウシアンスプラットの点光源可視性と光輸送をライト空間で定式化する。光源からの一回のラスタライズで沈着を記録し、受光点はそこからフィルタ付きの学習可能なシャドウ判定と、沈着した特徴に対して線形な光輸送項を導く。
\item 放射量に基づく段階的最適化により、細く入り組んだ幾何形状でもモデルを学習可能にする。対照実験を通じて、損失を計算する領域がそのような幾何形状の形成可否を決めることを示す。
\item 18 シーンの分析により、ライト空間の光輸送が有効な場合と局所ネットワークで十分な場合を明らかにする。また、現在の OLAT ベンチマークのテスト光源は学習光源からの角度差の中央値がわずか $1.1^\circ$ であり、光源の外挿をほとんど評価していないことを示す。
\end{itemize}
